\section{Week 7: Human-Centred Values -- Automation, Agency \& Fairness}

\subsection{Topic Summary}

\begin{keybox}[Learning Objectives]
By the end of this week, you will be able to:
\begin{itemize}
  \item \textbf{HAI3}: Apply critical reasoning skills to identify general (explainability, fairness, efficacy, accountability, transparency) and context-specific constraints on human-AI interactions.
\end{itemize}
\end{keybox}

\begin{keybox}[Big Picture]
This week begins the ``HAI as Values'' section---designing AI that is \textbf{aligned with human values}.

We cover three core values:
\begin{enumerate}
  \item \textbf{Automation}: When should AI act autonomously? When should humans retain control?
  \item \textbf{Agency}: How do we distribute responsibility between humans and AI?
  \item \textbf{Fairness}: How do we prevent AI from perpetuating or amplifying bias?
\end{enumerate}
\end{keybox}

\begin{keybox}[What Examiners Like to Ask]
\begin{itemize}
  \item Explain the shift from 1D to 2D thinking about automation and human control.
  \item Describe the Google Flu Trends and Boeing MCAS failures; what lessons do they teach?
  \item Define ``moral crumple zones'' and explain with an example.
  \item Distinguish between allocative and representational harms.
  \item Explain the Suresh \& Guttag framework for sources of bias in ML.
  \item Given a scenario, identify what type of bias might be present.
  \item Compare transparency and explainability.
\end{itemize}
\end{keybox}

\subsection{Overview + Core Concepts + Extended Theory}

\subsubsection*{Automation \& Agency}

\paragraph*{The Case for Automation}
Humans have limitations that AI can address:
\begin{itemize}
  \item Poor at repetitive tasks.
  \item Cognitive biases and judgment errors.
  \item Limited memory.
  \item Ambiguous intents.
\end{itemize}

\paragraph*{The Case Against Pure Automation}
\begin{itemize}
  \item Loss of critical domain expertise.
  \item Fairness concerns (who is harmed when AI fails?).
  \item Accountability gaps (who is responsible?).
\end{itemize}

\paragraph*{From 1D to 2D Thinking}

\textbf{Old (1D) Model}: Choose a point on a spectrum from human control to full automation.

\begin{center}
Human Control $\longleftrightarrow$ Full Automation
\end{center}

\textbf{New (2D) Model}: Aim for \textbf{both} high human control \textbf{and} high automation.

\begin{center}
\renewcommand{\arraystretch}{1.3}
\begin{tabular}{@{} l | c c @{}}
& \textbf{Low Automation} & \textbf{High Automation} \\
\midrule
\textbf{High Human Control} & Bicycle, piano & Camera (auto-focus + human click) \\
\textbf{Low Human Control} & --- & Pacemaker, airbag \\
\end{tabular}
\end{center}

Goal: Safe, reliable, trustworthy interactions where humans retain meaningful control \textit{and} benefit from automation.

\subsubsection*{Case Studies in Automation Failure}

\begin{exbox}[Google Flu Trends (2008--2015)]
\textbf{What it did}: Aggregated search queries (e.g., ``flu symptoms'') to predict flu outbreaks.

\textbf{The failure}: Predicted \textbf{more than double} the actual flu cases reported by CDC.

\textbf{Why it failed}:
\begin{itemize}
  \item Non-transparent algorithms.
  \item Features like ``autosuggest'' altered search patterns.
  \item Modeling became less effective over time.
\end{itemize}

\textbf{Lesson}: Automation without transparency and human oversight can fail silently.
\end{exbox}

\begin{exbox}[Boeing MCAS (2018--2019)]
\textbf{What happened}: Two crashes killed 346 people.

\textbf{The failure}: MCAS flight control software received incorrect sensor data (thought nose was up), continuously pushed nose down.

\textbf{Why it failed}:
\begin{itemize}
  \item Pilots couldn't override the automated system.
  \item Pilots weren't informed MCAS existed---no training.
\end{itemize}

\textbf{Lesson}: Automation failure + lack of human agency + lack of transparency = catastrophe.
\end{exbox}

\subsubsection*{Moral Crumple Zones}

\begin{defbox}[Moral Crumple Zone (Madeleine Elish)]
When responsibility for an action is \textbf{misattributed to a human} who had \textbf{limited control} over an automated system.

Like a car's crumple zone absorbs crash impact, the human ``absorbs'' moral/legal blame for AI failures.
\end{defbox}

\begin{exbox}[Uber Self-Driving Car (2018)]
A self-driving Uber struck and killed a pedestrian in Arizona. A ``safety driver'' was at the wheel but wasn't actively monitoring.

\textbf{Result}: The safety driver faced manslaughter charges---despite having limited control over the autonomous system.

\textbf{Question}: Who is really responsible when AI fails?
\end{exbox}

\subsubsection*{GenAI and Human Decision-Making}

\paragraph*{Study 1: ``To Use or Not to Use'' (CHI 2025)}
\textbf{Research questions}:
\begin{enumerate}
  \item How good are people at deciding when to use GenAI to maximize productivity?
  \item How do people decide when to use or not use GenAI?
\end{enumerate}

\textbf{Method}: Participants filled in image grids manually or with ``Assisted Fill'' (simulated GenAI). Varied latency and error rate.

\textbf{Findings}:
\begin{center}
\renewcommand{\arraystretch}{1.3}
\begin{tabular}{@{} l p{8cm} @{}}
\toprule
\textbf{AI Type} & \textbf{User Behavior} \\
\midrule
``Smart AI'' (fast + accurate) & Users make optimal decisions (use it). \\
``Dumb AI'' (slow + inaccurate) & Users make optimal decisions (don't use it). \\
Moderate AI & Users struggle to decide optimally. \\
Smart but slow & \textbf{Gulf of impatience}---users abandon too early. \\
Dumb but fast & \textbf{Gulf of overreliance}---users trust it too much. \\
\bottomrule
\end{tabular}
\end{center}

\paragraph*{Study 2: GenAI and Critical Thinking (CHI 2025)}
\textbf{Method}: Survey of 319 knowledge workers about GenAI use.

\textbf{Findings}:
\begin{itemize}
  \item Higher confidence in GenAI $\to$ \textbf{less critical thinking} (even though effort is low).
  \item Higher self-confidence $\to$ \textbf{more critical thinking} (even though effort is high).
\end{itemize}

\textbf{Design implications}:
\begin{itemize}
  \item AI tools should help users gauge reliability of outputs.
  \item Enhance awareness of when critical thinking is needed.
  \item Indicate when to trust AI vs.\ when to apply own judgment.
\end{itemize}

\subsubsection*{Fairness}

\begin{defbox}[Fairness]
The \textbf{absence of prejudice or favoritism} toward an individual or group based on their inherent or acquired characteristics.

An \textbf{unfair algorithm} is one whose decisions are skewed toward a particular group.
\end{defbox}

\paragraph*{Why Fairness Matters}
AI is used in high-stakes decisions:
\begin{itemize}
  \item Courts (recidivism prediction).
  \item Healthcare.
  \item Child welfare.
  \item Autonomous vehicles.
  \item Hiring and loans.
\end{itemize}

\begin{exbox}[COMPAS Algorithm (ProPublica, 2016)]
\textbf{What it did}: Predicted likelihood of defendants re-offending; used by US courts for bail/detention decisions.

\textbf{The bias}: Higher false positive rates for African-Americans (predicted high risk when they didn't re-offend).

\textbf{Key finding}: COMPAS was no better than non-expert human judgment---but with racial bias built in.

\textbf{Lesson}: Automated decision-making can encode and amplify societal biases.
\end{exbox}

\begin{exbox}[Gender Shades (Joy Buolamwini)]
\textbf{What she found}: Commercial facial recognition had:
\begin{itemize}
  \item 99\% accuracy for white men.
  \item 35\% accuracy for dark-skinned women.
\end{itemize}

\textbf{The ``coded gaze''}: Algorithms trained on biased data perpetuate that bias.

\textbf{Impact}: Led to major tech companies improving their systems and policy discussions.
\end{exbox}

\subsubsection*{Types of Harm (Kate Crawford)}

\begin{center}
\renewcommand{\arraystretch}{1.3}
\begin{tabular}{@{} l p{5cm} p{5cm} @{}}
\toprule
& \textbf{Allocative Harms} & \textbf{Representational Harms} \\
\midrule
Definition & System withholds opportunity/resource from a group & System reinforces subordination of groups \\
Effect & Immediate (decision-making) & Long-term (culture, identity) \\
Example & Denied loan, detained by algorithm & Stereotypes in search results \\
Formalizable & Easier (resource allocation) & Harder (cultural impact) \\
Research focus & More attention historically & Less attention, but critical \\
\bottomrule
\end{tabular}
\end{center}

\subsubsection*{Sources of Bias in ML (Suresh \& Guttag Framework)}

\begin{center}
\renewcommand{\arraystretch}{1.4}
\begin{tabular}{@{} l p{5cm} p{5cm} @{}}
\toprule
\textbf{Bias Type} & \textbf{Description} & \textbf{Example} \\
\midrule
Historical & Data reflects real-world biases, even if perfectly measured & Word embeddings associate ``nurse'' with women, ``engineer'' with men \\
Representation & Sample doesn't represent full population & ImageNet: 1--2\% images from China/India; fails on non-Western weddings \\
Measurement & Proxy features misrepresent underlying concept & COMPAS uses ``arrest'' as proxy for ``riskiness''---but minorities are over-policed \\
Learning & Model choices amplify disparities & Pruning for compactness preserves frequent features, loses minority attributes \\
\bottomrule
\end{tabular}
\end{center}

\paragraph*{Historical Bias}
Even if data is perfectly sampled, it can perpetuate harmful stereotypes because those stereotypes exist in the world.

\paragraph*{Representation Bias}
\begin{itemize}
  \item Target population doesn't reflect reality.
  \item Underrepresented groups (e.g., pregnant women in 18--60 sample).
  \item Uneven geographic sampling (ImageNet: mostly US/Western Europe).
\end{itemize}

\paragraph*{Measurement Bias}
Using proxies that don't accurately measure the concept:
\begin{itemize}
  \item ``Arrest'' $\neq$ ``riskiness'' (arrests are influenced by policing patterns).
  \item Minority communities are more heavily policed $\to$ more arrests $\to$ higher ``risk'' scores.
\end{itemize}

\paragraph*{Learning Bias}
Model optimization can amplify disparities:
\begin{itemize}
  \item Pruning for compact models prioritizes frequent features.
  \item Underrepresented attributes are lost.
\end{itemize}

\subsubsection*{Transparency \& Explainability}

\begin{defbox}[Transparency]
AI systems should not operate as ``black boxes.'' Users should have insight into the AI's reasoning process.
\end{defbox}

\begin{defbox}[Explainability]
The ability of AI systems to \textbf{justify their outputs} in a way humans can understand.
\end{defbox}

\begin{center}
\renewcommand{\arraystretch}{1.3}
\begin{tabular}{@{} l p{5.5cm} p{5.5cm} @{}}
\toprule
& \textbf{Transparency} & \textbf{Explainability} \\
\midrule
Focus & How the system works & Why a specific output was produced \\
Audience & Developers, auditors & End users \\
Goal & Enable oversight & Enable trust and correction \\
\bottomrule
\end{tabular}
\end{center}

Both are mechanisms for addressing fairness concerns by making bias visible and correctable.

\subsubsection*{Key Takeaways}

\begin{enumerate}
  \item \textbf{Automation without agency is dangerous}: Users need ability to override, and training to do so.
  \item \textbf{Transparency matters}: Non-transparent systems fail silently (Google Flu Trends).
  \item \textbf{Responsibility is distributed}: Moral crumple zones misattribute blame to humans with limited control.
  \item \textbf{GenAI changes behavior}: Overreliance and reduced critical thinking are real risks.
  \item \textbf{Fairness is multi-dimensional}: Allocative and representational harms; historical, representation, measurement, learning bias.
  \item \textbf{Design choices matter}: From data collection to model architecture, every step can introduce bias.
\end{enumerate}

\subsubsection*{Key Definitions Summary}

\begin{center}
\renewcommand{\arraystretch}{1.3}
\begin{tabular}{@{} l p{8cm} @{}}
\toprule
\textbf{Term} & \textbf{Definition} \\
\midrule
Fairness & Absence of prejudice/favoritism based on characteristics \\
Allocative harm & Withholding opportunity/resource from a group \\
Representational harm & Reinforcing subordination/stereotypes \\
Historical bias & Data reflects real-world biases \\
Representation bias & Sample doesn't represent full population \\
Measurement bias & Proxy features misrepresent underlying concept \\
Learning bias & Model optimization amplifies disparities \\
Moral crumple zone & Human absorbs blame for AI failure despite limited control \\
Transparency & Insight into AI's reasoning process \\
Explainability & AI justifies outputs in human-understandable way \\
\bottomrule
\end{tabular}
\end{center}

\subsubsection*{Recommended Reading}
\begin{itemize}
  \item \textbf{Textbook}: Shneiderman, Human-Centered AI, Chapter 8
  \item \textbf{Video}: Kate Crawford, ``The Trouble with Bias'' (NeurIPS 2017)
  \item \textbf{Project}: Joy Buolamwini, Gender Shades
  \item \textbf{Article}: ProPublica, ``Machine Bias'' (COMPAS investigation)
  \item \textbf{Paper}: Suresh \& Guttag, ``Framework for Understanding Sources of Harm''
\end{itemize}


